\chapter{推荐概述}
\label{chap:rs-overview}

准备一次登山旅行时，一个人可能先在首页看到背包和水壶，随后搜索“防水背包，500元以内”，又在结果中遇到商家的赞助商品。屏幕上都是商品，系统承担的任务却已经改变：首页需要从有限历史推断哪些东西值得提出，搜索需要兑现用户刚刚写出的条件，广告还要决定哪些商家有资格竞争、怎样收费以及怎样使用剩余预算。解释这三种应用，不能只问采用了哪个预测模型，还要追踪需求依据怎样成为候选、候选怎样成为展示，以及展示留下的反馈能够证明什么。

前面的模型与范式提供了表示、预测和学习工具。本部分把这些工具放回实际选择过程，沿理解、选择与展示建立共同任务，再讨论多方目标、连续决策与效果证据。搜索和广告在这一共同基础上展开各自的完整任务。读完本章，应当能够指出一个方法改变了哪些输入、输出或学习关系，并知道它尚未承担哪些职责。

% Outline: 1.1
\section{推荐的应用}
\label{sec:rs-applications}

当目录只有十件商品时，用户可以逐一查看；当目录包含大量对象而注意时间基本不变时，选择就需要帮助。简单按销量排列能够利用群体经验，却无法同时服务刚开始登山的人、准备更换轻量装备的人和只想寻找一个通勤包的人。\term{推荐系统}（recommender system）根据当前可用的需求与对象信息，选择值得向用户提供的对象或交互。个性化表示选择可以随人及请求条件变化，不表示每次都必须识别一个长期稳定的用户身份。

推荐并非只发生在用户不知道要什么的时候。用户可能知道用途，却不知道有哪些产品；也可能已经选好相机，需要比较适配的配件。系统提供的帮助相应包括发现、比较、补充和避免无效选择。热门先验对新用户有用，明确的预算和用途对当前请求更直接。历史偏好只有在与本次任务相容时，才是有价值的补充依据。

相同的输入格式可以对应不同的消费过程。看一段短视频通常很快产生观看或跳过记录，购买一个背包可能需要比较数日，预约本地服务还受到地点、营业时间与容量限制。表~\ref{tab:rs-scenarios}按这些条件比较场景。表中没有哪种场景天然更容易：反馈密集可能同时包含大量偶然点击，反馈稀疏则可能更接近有成本的选择。

\begin{table}[htbp]
  \centering
  \begin{tabularx}{\linewidth}{@{}lXX@{}}
    \toprule
    场景 & 影响选择的条件 & 反馈解释的难点 \\
    \midrule
    内容消费 & 注意有限，部分内容寿命短 & 观看时长受内容长度和退出成本影响 \\
    电商商品 & 价格、规格及兼容性共同作用 & 购买稀疏，成交可能晚于首次浏览 \\
    本地服务 & 时间、距离及容量限制可用性 & 点击不代表可以到店或实际完成服务 \\
    专业信息 & 正确性、来源与证据完整性重要 & 打开文档不代表问题已经解决 \\
    \bottomrule
  \end{tabularx}
  \caption{应用条件决定选择任务和反馈含义。相同的点击事件不能抹去消费成本、结果时延及对象可用性的差别。}
  \label{tab:rs-scenarios}
\end{table}

在登山旅行的例子中，首页看到用户近期浏览了背包、水壶和相机，可以把户外与摄影作为两类待检验的需求依据。用户搜索“防水通勤背包，价格不超过500元”之后，当前查询给出了更明确的方向：相机历史不应把结果带离背包，过去购买过高价商品也不应覆盖价格上限。赞助展示继续使用这些条件，但参与者增加了广告主。一个广告是否相关、是否合格、预期会产生多少价值、最后支付多少钱，是互相关联而又不同的判断。

选择的单位也不总是单个对象。三个单独得分很高的背包可能非常相似，组成一屏后并没有帮助用户比较。一个家庭共同选择电影时，个人平均得分可能掩盖某位成员强烈反对的情况。此时需要把列表或群组请求作为任务对象，保留成员、对象关系和展示位置等信息。社交、时间和地点可以补充判断依据；隐私限制与跨域身份缺失又会减少可用信息。后续方法必须在这些具体条件下定义，不能把它们压成同一种“用户向量不足”的问题。

% Outline: 1.2
\section{推荐的基本概念}
\label{sec:rs-concepts}

为了把应用过程写清楚，先固定一次选择发生在什么时点。以$u\in U$表示用户，$i\in I$表示目录中的物品；“物品”在此泛指可选择的真实条目，可以是商品、视频、文档或服务。用户身份连接被允许使用的信息，物品标识连接对象事实和状态。一次\term{请求}（request）则给出当前决策时点$t$，以及这次选择的入口、表达和可观察条件。同一用户在首页和搜索页发出的请求，可能具有不同目标。

\term{情境}（context）包括时间、设备、地点、会话状态等当前条件。历史$H_{u,t}$由请求之前已经可用的事件组成；每条事件至少需要说明对象、行为类型和发生时间，涉及学习时还要保存当时的展示机会及记录条件。发生在今天的购买可以作为昨天预测的训练目标，但不能倒过来进入昨天请求的输入。图~\ref{fig:rs-request-evidence}把当前请求的证据与之后发生的反馈分开。

\begin{figure}[htbp]
  \centering
  % 请求前证据与请求后反馈；最终宽度约109mm。
\begin{tikzpicture}[
  font=\figurefont\footnotesize,
  rsbox/.style={rectangle,draw=BookRule,fill=BookPaper,
    text=BookInk,align=center,text width=28mm,minimum height=12mm,inner sep=3pt},
  rsflow/.style={-{Stealth[length=1.8mm]},draw=BookMuted,line width=.8pt}
]
  \node[rsbox] (history) at (0,1.6) {已到达的历史\\浏览、收藏、购买};
  \node[rsbox] (request) at (3.8,1.6) {当前表达与情境\\用途、价格、时点};
  \node[rsbox] (items) at (7.6,1.6) {对象资料与状态\\属性、库存、版本};
  \node[rsbox,fill=FigureModel!10] (decision) at (3.8,0) {本次决策\\选择与展示};
  \node[rsbox,fill=FigureLoss!10] (feedback) at (3.8,-1.6) {之后收到的反馈\\供后续请求使用};
  \coordinate (join) at (3.8,.8);
  \draw[draw=BookMuted,line width=.8pt] (history.south) |- (join);
  \draw[draw=BookMuted,line width=.8pt] (items.south) |- (join);
  \draw[rsflow] (request.south) -- (decision.north);
  \draw[rsflow] (decision.south) -- (feedback.north);
  \draw[rsflow,dashed,rounded corners=1mm]
    (feedback.west) -- (-1.9,-1.6) -- (-1.9,1.6) -- (history.west);
\end{tikzpicture}

  \caption{一次请求使用的证据与事后反馈。虚线表示反馈在到达后才可用于更新；它不能返回成为同一次已执行决策的输入。}
  \label{fig:rs-request-evidence}
\end{figure}

用户明确打分、声明偏好或拒绝某类内容，属于\term{显式反馈}（explicit feedback）：记录直接表达了某种判断，但仍需保留问法、对象和时间。点击、观看与购买属于\term{隐式反馈}（implicit feedback）：系统观察到了行为，还需要结合机会与情境解释其含义。没有点击可能因为不感兴趣，也可能因为没有看到；完成购买也可能反映原有需求，而非推荐带来的额外作用。

匹配依据还可以按来源区分。\term{内容依据}（content-based evidence）来自物品属性、文本、图片及其他表达；\term{协同依据}（collaborative evidence）来自不同用户与物品之间的交互关系。例如，用户经常共同购买相机和存储卡是一种交互关系，商品资料明确写出接口兼容则是另一种证据。前者不能自动证明后者，二者却可以结合使用。一个模型使用点击标签，只能说明监督来自行为；是否利用协同关系还要检查它怎样共享用户、物品或交互信息。

选择过程需要区分几个输出。\term{候选召回}（candidate retrieval）从当前可用目录$I_t$中形成候选集合$C_t$，决定后续有机会比较哪些物品。\term{响应预估}（response prediction）估计在给定条件下某种行为发生的概率或结果数值。\term{排序}（ranking）确定候选之间的次序，可以学习相对偏好，也可以使用响应或价值估计。\term{列表构造}（slate construction）确定共同展示的对象及顺序；\term{混排}（blending）进一步组合自然内容、广告等不同来源。这些任务能够联合实现，但它们要求的输出仍需明确。

用$s_\theta(\bm{x}_t,i)$表示参数为$\theta$的模型给候选$i$的分数，其中$\bm{x}_t$汇集本次允许使用的请求特征。这个记号本身没有规定分数含义。同一个数值0.8，可以表示模型对偏好比较的评分，也可以表示经过特定训练得到的点击概率。只有后者满足概率的定义及使用条件时，才可写作$p_\theta(y=1\mid\bm{x}_t,i)$并据此计算期望响应。

下面四种量尤需分清。\termalias{search-relevance}{相关性}（relevance）描述对象在多大程度上满足当前需求；响应概率描述给定条件下行为发生的可能性；消费价值说明行为或结果有什么用途；\term{因果增量}（causal effect）比较采用与不采用某项行动会造成多大差异。高点击概率的标题可能误导读者，高购买概率的商品可能本来就会被购买。因而“分数更高”只有在目标量明确后才具有可比较的含义。

\begin{example}[高响应与高增量可以指向不同对象]
  \label{ex:rs-response-effect}
  假设在同一种机会和结果窗口下，商品A展示与不展示时的购买概率分别为0.90和0.85，商品B分别为0.60和0.30。这组概率是教学假设，不是由普通点击日志直接得到的事实。按展示后的购买概率选A，按概率差比较则得到A的增量0.05、B的增量0.30，从而选B。两种选择回答不同问题；实际取得这四个量还需要对照证据，见第~\ref{chap:rs-evaluation}~章。
\end{example}

另外几组名称相近的概念位于不同层面。浏览、收藏和购买一起进入模型，是多行为输入；模型同时预测点击和购买，是多任务输出；平台在收入与体验之间作出取舍，是多目标决策。输入更多行为，不意味着各行为可以共用标签；预测多个结果，也没有自动确定它们之间的价值权重。

\term{概率校准}（probability calibration）检验预测为0.1的一组机会，成熟结果的发生频率是否接近0.1。\term{列表校准}（calibrated recommendation）在本部分指列表的类别构成与某种有依据的兴趣构成相对应。预测不确定性则描述证据对判断的支持程度；控制推荐集合的平均错误风险还需要额外的校准样本和统计条件。这些要求分别作用于概率、构成、证据和集合风险，不能用“已经校准”一语替代。

数据使用也有两个时间方向。隐私保护限制从训练信息及输出中能够推断什么，数据撤回要求处理此前已经学入模型的影响。记录留在设备上不自动证明更新没有泄漏，删除一条日志也不自动消除其对参数或协同表示的作用。第~\ref{chap:rs-ranking}~章和第~\ref{chap:rs-evaluation}~章将分别说明学习、维护与验收的接口。

\subsection{推荐系统的参考架构}
\label{sec:rs-reference-architecture-concept}

有了这些对象，可以沿一次请求追踪各任务的关系。理解结合用户表达、历史与物品事实，形成需求和匹配依据；候选召回把这些依据用于当前目录；排序或响应预估为候选提供比较信息；展示决定用户真正有机会看到和操作什么。图~\ref{fig:rs-reference-architecture}中的箭头表示输入输出依赖，反馈则在之后更新信息、模型与策略。

\begin{figure*}[htbp]
  \centering
  % 实线为请求中的数据，虚线为事后更新；最终宽度约164mm。
\begin{tikzpicture}[
  font=\figurefont\footnotesize,
  rsbox/.style={rectangle,draw=BookRule,fill=BookPaper,
    align=center,text=BookInk,text width=25mm,minimum height=13mm,inner sep=3pt},
  rsflow/.style={-{Stealth[length=1.8mm]},draw=BookMuted,line width=.9pt},
  rslabel/.style={font=\figurefont\footnotesize,text=BookInk,fill=white,inner sep=1pt}
]
  \node[rsbox,fill=FigureInput!12] (input) at (0,0) {请求与目录\\历史、条件、事实};
  \node[rsbox,fill=FigureModel!12] (understanding) at (3.35,0) {理解\\需求与对象依据};
  \node[rsbox,fill=FigureModel!12] (recall) at (6.7,0) {召回\\候选集合};
  \node[rsbox,fill=FigureModel!12] (rank) at (10.05,0) {预估与排序\\候选比较信息};
  \node[rsbox,fill=FigureOutput!10] (display) at (13.4,0) {展示与交互\\实际机会与结果};
  \draw[rsflow] (input.east) -- (understanding.west);
  \draw[rsflow] (understanding.east) -- (recall.west);
  \draw[rsflow] (recall.east) -- (rank.west);
  \draw[rsflow] (rank.east) -- (display.west);
  \node[rsbox,text width=56mm,fill=FigureLoss!10] (state) at (8.35,1.9)
    {累计状态与要求\\预算、曝光、频次及多方目标};
  \draw[rsflow] (state.south) -- (8.35,1) -| (recall.north);
  \draw[rsflow] (8.35,1) -| (rank.north);
  \draw[rsflow] (8.35,1) -| (display.north);
  \node[rsbox,text width=43mm,fill=FigureLoss!10] (learning) at (8.35,-2.15)
    {反馈学习与效果检验\\标签、模型与策略更新};
  \draw[rsflow,rounded corners=1mm]
    (display.south) |- (learning.east);
  \draw[rsflow,dashed,rounded corners=1mm]
    (learning.west) -| (understanding.south);
\end{tikzpicture}

  \caption{推荐、搜索和广告的共同参考架构。实线追踪当前请求的选择过程，虚线表示事后更新。累计状态可影响多个任务；图不规定拍卖与混排必然采用某一种执行顺序。}
  \label{fig:rs-reference-architecture}
\end{figure*}

这一架构划分的是职责，不是要求所有系统串行调用四个独立模型。用户理解和物品理解可以并行，召回和排序可以共享用户表示，一个生成模型也可以直接输出物品列表。判断职责是否已被承担，要看实际输出：输出用户向量尚未取得真实物品，输出物品标识尚未确认库存，输出点击概率尚未决定整页负载，输出广告次序尚未定义支付。

以候选生成与排序的分工为例，前者在很大目录上追求足够覆盖，后者在较小集合上使用更精细的信息。YouTube在2016年报告的两阶段方案用用户网络和物品嵌入形成候选，再由独立网络排序；候选模型与排序模型的监督和输出用途并不相同。这一实例说明阶段分工可以来自目录规模、目标和计算预算，而不能被当作所有应用必须遵循的固定模板\scite{rs-r01}。

其中的计算分工可以沿一次观看请求理解。候选网络将已知观看和搜索等信息汇成用户向量$\bm h_t$，每个视频有可提前保存的向量$\bm v_i$，两者的内积$\bm h_t^\top\bm v_i$提供访问目录的依据。学习时用实际的下一观看视频作为目标，让它在训练竞争集合中得到更高概率；目录很大时采用采样近似。服务时只计算一次$\bm h_t$，再从物品向量索引取回候选。第~\ref{sec:rs-youtube-candidate}~节给出目标函数、采样与预测时点的具体条件。

排序网络接收已经取回的候选，可以进一步使用当前候选与历史、情境的交互特征。它需要为较小候选池逐项计算，而不必对整个目录使用这套较贵的计算。学习候选访问的表示与学习展示后的响应，因而是两个可以分工的任务。实际应用还要检查候选覆盖、曝光反馈和服务时延；一个阶段更准确，不能自动补上另一个阶段遗漏的对象。

工业系统常把选择继续分成\term{粗排}（pre-ranking）与\term{精排}（fine-ranking）。粗排先用较便宜的评分缩小较大候选池，精排再为剩余候选使用更丰富的输入和计算。例如，从5000个候选选500个、再选50个，是一个教学配置，数量由服务预算和质量要求决定。若某个相关对象在粗排被截去，精排仍无法找回它；因此，阶段之间除了共享表示，还需要核验目标与截断是否相容。

候选截断还产生一个简单而重要的边界。若最终列表$L_t$只从$C_t$中选择，则
\begin{equation}
  \{i:i\in L_t\}\subseteq C_t\subseteq I_t.
  \label{eq:rs-candidate-containment}
\end{equation}
因此，排序器无论多精细，都不能选出候选池之外的对象。需要查明错误发生在需求误读、候选遗漏还是候选比较，才知道应增加信息、调整召回还是改进评分。直接目录列表生成改变了中间候选接口，但同样需要说明它实际能够访问哪些物品，以及解码预算是否造成遗漏。

反馈进入多个任务，是因为同一事件可以回答不同问题。一次“太重了”的拒绝为理解补充重量条件，可以帮助候选排除不合格对象，也能作为展示交互是否兑现要求的检查。如果把它用于训练，还需要知道拒绝针对哪个对象和哪个表达；如果用它证明新策略更好，还需要适当的比较设计。取得反馈、解释反馈和识别策略效果，不能因共用日志就被合并。

% Outline: 1.3
\section{推荐的研究内容}
\label{sec:rs-research-problems}

沿上一节的参考架构，可以把推荐面对的研究问题按一次请求的处理链拆开。每个问题都有自己要改变的量：理解改变可用证据的组织，召回改变能够访问哪些对象，排序改变候选之间的比较，展示改变用户实际面对的整组内容；再向外延伸，多方目标、连续决策和效果证据分别改变优化的取舍、时间范围和判断依据。表~\ref{tab:rs-research-map}把这些问题、主要挑战与后续章节对应起来。

\begin{table}[htbp]
  \centering
  \begin{tabularx}{\linewidth}{@{}lXl@{}}
    \toprule
    研究问题 & 关键挑战 & 对应章节 \\
    \midrule
    需求与对象怎样理解 & 隐式反馈的机会解释、行为顺序与长历史、稀疏对象与新对象 & 第~\ref{chap:rs-understanding}~章 \\
    候选怎样取得 & 大目录下的覆盖与预算、可预计算与实时性、截断不可恢复、动态目录 & 第~\ref{chap:rs-retrieval}~章 \\
    候选怎样比较 & 目标量的区分、特征组合、多行为多任务、校准与选择偏差 & 第~\ref{chap:rs-ranking}~章 \\
    结果怎样展示 & 列表与群组决策、多样与重复、位置表达、混排与探索 & 第~\ref{chap:rs-presentation}~章 \\
    多方目标怎样取舍 & 收入与体验等多目标、资源与资格约束、价值权重 & 第~\ref{chap:rs-constraints}~章 \\
    长期影响怎样控制 & 反馈循环、兴趣与供给随曝光变化、长期与即时价值 & 第~\ref{chap:rs-long-term}~章 \\
    效果由什么证据支持 & 离线与在线差距、因果增量、标签成熟、强基线与实验 & 第~\ref{chap:rs-evaluation}~章 \\
    显式查询怎样满足 & 查询理解、相关性、稀疏与稠密检索、会话修正 & 第~\ref{chap:rs-search}~章 \\
    付费投放怎样决策 & 资格、响应链路、竞价、跨请求预算与增量测量 & 第~\ref{chap:rs-advertising}~章 \\
    \bottomrule
  \end{tabularx}
  \caption{推荐的研究问题沿处理链分布，各问题面对不同挑战并由后续章节展开。同一事件或模型可以同时涉及几行，划分依据是它要改变的决策量。}
  \label{tab:rs-research-map}
\end{table}

第一类问题是需求与对象怎样理解。显式反馈稀少且带问法，隐式反馈又混有位置、机会和偶然因素；行为有顺序、有时延，完整历史越长，相关事件的访问和压缩越困难。稀疏对象和新对象缺少交互证据，只能更多依靠内容依据或跨对象迁移；群组、跨域和隐私限制还会改变哪些主体与信息可以对应。这些问题由第~\ref{chap:rs-understanding}~章处理。

第二类问题是候选怎样从大目录中取得。目录规模使模型不可能为每个对象逐项使用最重的计算，必须在覆盖、时延和预算之间安排可预计算与实时计算；式~\eqref{eq:rs-candidate-containment}的包含关系又意味着召回遗漏无法靠排序补救。动态目录还要求新物品进入、失效物品退出，而不是只在静态测试集上报告命中。这对应第~\ref{chap:rs-retrieval}~章。

第三类问题是候选之间怎样比较。排序先要分清自己估计的是相关性、响应概率、消费价值还是因果增量，再决定特征和历史怎样组合、多个行为与多个目标怎样共同学习。点击日志还带有选择偏差和位置偏差，概率分数是否可作为期望值使用取决于校准与样本条件。这对应第~\ref{chap:rs-ranking}~章。

第四类问题是结果怎样组织成一次展示。决策对象常常是一屏列表或一个群组而非单个物品，需要在准确匹配之外处理多样性、重复、位置、表达方式以及自然内容与广告的混排；让新对象取得曝光的探索还要权衡信息收益与当次体验。这对应第~\ref{chap:rs-presentation}~章。

处理链之外还有三组贯穿性问题。平台需要在收入、体验和提供者生态等多个目标，以及资源与资格约束之间作出取舍，价值权重不能由预测模型自动给出，这是第~\ref{chap:rs-constraints}~章；旧策略决定了下一轮能观察到哪些反馈，持续曝光又可能改变兴趣和供给，单次响应好不等于长期价值好，这是第~\ref{chap:rs-long-term}~章；而上述任何判断是否成立，都要由反事实推断、在线实验和可靠的离线比较提供证据，这是第~\ref{chap:rs-evaluation}~章。搜索把显式查询和相关性贯穿处理链，广告在资格和价值之外增加竞价、支付与跨请求预算，分别由第~\ref{chap:rs-search}~章和第~\ref{chap:rs-advertising}~章展开。

% Outline: 1.4
\section{推荐的研究历史}
\label{sec:rs-development}

推荐技术的发展可以按几个阶段理解，但阶段之间不是替代关系：许多长期问题从早期延续至今，新的网络结构并没有让它们消失。

最早的系统化方案围绕协同过滤展开。1994年的GroupLens以新闻评分建立协同过滤系统：过去对文章判断相近的人，为彼此尚未读过的文章提供预测依据。2001年的物品协同过滤把可复用的相似关系移到物品之间，使请求时可以从用户已交互对象的邻居生成预测；这里的转折包括计算怎样复用，并非只把用户相似改成物品相似\scite{rs-h07,rs-h03}。与协同路线并行，直接属性匹配使用物品已有描述。矩阵分解则通过共享潜在因子联系稀疏观测，潜在因子是从记录中学习的少量数值，不要求逐维对应人工命名的兴趣类别；2009年的矩阵分解综述总结了这种路线在推荐任务中的目标、偏置与扩展方式\scite{rs-h08}。内容、邻域协同与潜在因子依据不同，却一直可以组合使用；后来的深层网络也没有让稀疏观测、未曝光对象和新物品的证据问题自然消失。

第二阶段关注特征之间的组合怎样计算。我们把用户类别、商品类别、设备等一组组输入称为\term{字段}（field）；字段中的具体取值是特征。例如，“设备”是一个字段，“手机”是其中一个取值。线性模型为每个特征赋参数，显式交叉为“手机与视频”“通勤用途与背包”等组合另设参数；组合很多且观察稀疏时，逐对参数不容易取得足够证据。2010年的因子分解机（factorization machine, FM）为特征学习短向量，以两向量的内积表示交互，同一个向量可以参与多种组合，从而共享学习依据。2016年的字段感知因子分解机（field-aware factorization machine, FFM）进一步让一个特征根据对方所属字段使用不同向量，例如设备与商品、设备与位置的交互分别表示。两者的计算、学习目标和参数代价在第~\ref{sec:rs-fm}~节与第~\ref{sec:rs-ffm}~节展开\scite{rs-h05,rs-h06}。此后深层交互、字段注意和容量扩展继续研究哪些组合可以表达、计算如何分配，但改变的始终是候选评分内部的计算，候选输入和分数输出可以保持不变。

与特征组合并行，行为建模追问历史中的哪些信息值得保留。循环状态、因果自注意、掩码重建和会话图使用不同的顺序与监督结构；DIN按候选选择历史，长历史方法进一步比较固定容量记忆、检索相关事件、哈希聚合和层次压缩。完整历史越长，模型获得的信息可能越丰富，访问、训练和更新成本也随之变化；比较时要同时控制可用事件、候选数量和时延，不能把“最长可读多少条历史”当作理解质量本身。

目录规模和服务预算则推动了另一条沿候选访问展开的路线，把可提前计算的物品表示用于索引，或者沿树路径逐步访问物品。工业系统因此普遍形成候选、粗排、精排和列表决策的职责分担，YouTube两阶段方案是这一分工的早期代表。比较独立编码、候选相关历史匹配与长历史检索时，需要看清计算移到哪一侧、哪些结果可预先计算，以及相同历史怎样服务不同候选。语义标识生成进一步改变了访问方式，将历史映射为物品代码再查回真实目录；连续向量生成仍可以接近邻目录，交互分布重建也仍可以输出全目录分数。这里要按最终怎样取得对象来分类，而不能由训练过程出现“生成”二字就推断输出是物品标识。

近期的工作尝试把这些阶段重新整合，表~\ref{tab:rs-change-scope}区分了几种经常被统称为“统一”的变化。其中，评分骨干指从字段、历史等输入到响应头之前的主要计算网络；交互层只是其中的一部分。各行不是升级顺序，也不互斥：同一个方案可以同时改进字段交互、共享骨干和跨阶段训练。例如，阶段共享可以保留候选逐层缩减，列表生成则可能取消外部候选池，却仍需执行对象资格和展示规则。

\begin{table}[htbp]
  \centering
  \begin{tabularx}{\linewidth}{@{}lXX@{}}
    \toprule
    改变的范围 & 直接改变的关系 & 核验对象 \\
    \midrule
    字段交互 & 同一候选的特征怎样组合 & 分数、参数及计算量 \\
    评分骨干 & 字段与历史怎样共同计算 & 可见信息、共享参数及响应头 \\
    跨阶段共享 & 召回和各评分阶段复用什么 & 阶段目标、候选截断及梯度路径 \\
    标识生成 & 怎样从代码恢复物品 & 合法代码、目录映射及维护 \\
    列表生成 & 多个物品怎样共同产生 & 整组依赖、重复及执行资格 \\
    \bottomrule
  \end{tabularx}
  \caption{按实际改动范围理解架构演进。模型名称相近或采用同一种网络，不足以证明其输入、输出和整合范围相同。}
  \label{tab:rs-change-scope}
\end{table}

OneTrans将行为与字段放入共同评分堆栈，仍为给定候选产生响应分数。OneTrans-V2把共享扩展到召回、粗排和精排，同时保留阶段特征、候选缩减以及精排到粗排的训练蒸馏。共享骨干与逐阶段执行因而可以同时成立；这里采用其2026年9月的预印本，不把平台实验外推为所有推荐流程的替代结论\scite{rs-r43}。OneRec研究的是另一种整合范围：由历史直接生成多个物品的标识序列，再利用列表偏好反馈调整生成模型。模型可以在生成中学习对象之间的关系，但生成结果仍需恢复真实对象并满足展示条件；奖励模型偏好的列表、可执行的列表与用户实际受益的列表，仍然需要分别验证\scite{rs-r08}。

语言模型用于推荐的研究通常称为LLM4Rec，但这一名称描述的是工具及研究范围，并没有固定其应用角色。语言模型可以提取对象属性、整理需求，也可以给候选判定、生成物品标识、调用信息工具或产生解释。P5把评分、序列推荐、直接推荐、解释及评论等任务写成共享文本输入输出；任务提示和目标文本定义不同的学习问题，共用语言接口并未使它们获得相同的目录访问能力或效果证据\scite{rs-r21}。因此可以沿输出检查语言模型到底做了什么：输出属性时要核验对象事实，输出画像时要回到已有证据，输出代码时要查回目录，输出候选次序时要检查遗漏和重复，输出理由时要区分事实准确与对评分器忠实；调用信息工具时，还要把额外资料带来的改善与规划方式分开对照。图~\ref{fig:rs-interface-comparison}固定请求和目录，把外部候选评分、生成标识和生成列表三条路线的接口放在一起比较：它们可以共享表示和训练数据，但中间输出与截断位置不同，成本统计也必须覆盖实际执行的全过程。

\begin{figure}[htbp]
  \centering
  % 三种接口的平行对照，每条路径只表达对象交接。
\begin{tikzpicture}[
  font=\figurefont\footnotesize,
  rsbox/.style={rectangle,draw=BookRule,fill=BookPaper,
    align=center,text=BookInk,text width=24mm,minimum height=11mm,inner sep=3pt},
  rsflow/.style={-{Stealth[length=1.8mm]},draw=BookMuted,line width=.8pt},
  rshead/.style={font=\figurefont\bfseries\footnotesize,text=BookInk,anchor=west}
]
  \node[rshead] at (-1.3,0.9) {目录检索与候选评分};
  \node[rsbox] (r1) at (0,0) {目录访问};
  \node[rsbox] (r2) at (3.7,0) {候选集合\\逐候选评分};
  \node[rsbox] (r3) at (7.4,0) {列表组织\\核验与展示};
  \draw[rsflow] (r1.east) -- (r2.west);
  \draw[rsflow] (r2.east) -- (r3.west);
  \node[rshead] at (-1.3,-1.2) {物品标识生成};
  \node[rsbox] (s1) at (0,-2.1) {历史与条件\\生成物品代码};
  \node[rsbox] (s2) at (3.7,-2.1) {目录映射\\候选集合};
  \node[rsbox] (s3) at (7.4,-2.1) {评分及列表\\核验与展示};
  \draw[rsflow] (s1.east) -- (s2.west);
  \draw[rsflow] (s2.east) -- (s3.west);
  \node[rshead] at (-1.3,-3.3) {目录列表生成};
  \node[rsbox] (l1) at (0,-4.2) {历史与条件\\生成多物品代码};
  \node[rsbox] (l2) at (3.7,-4.2) {目录映射\\候选列表};
  \node[rsbox] (l3) at (7.4,-4.2) {列表核验\\执行展示};
  \draw[rsflow] (l1.east) -- (l2.west);
  \draw[rsflow] (l2.east) -- (l3.west);
\end{tikzpicture}

  \caption{三种选择接口的对照。框中列出各步骤的输入、处理与交接结果。生成标识可以保留后续评分，生成列表也仍须核验真实对象与展示规则；图不表示质量高低。}
  \label{fig:rs-interface-comparison}
\end{figure}

与结构演进相伴，效果证据的标准也在收紧。Ferrari Dacrema等人在2019年复查2015—2018年四个会议中的18项神经Top-$N$推荐研究，按其代码、数据和切分可恢复的标准复现了7项，再加入调优的邻域与图方法；6项没有在所考察配置中持续胜过这些简单基线，该范围排除了会话和群组推荐，不能外推成所有神经模型都无收益\scite{rs-e13}。一个具体对照是神经矩阵分解NeuMF：在原研究的Pinterest切分与命中协议中，复现实验的物品近邻HR@10为0.8744，NeuMF为0.8719；在MovieLens-1M中，NeuMF为0.7120，高于该物品近邻的0.6596，但调优的稀疏线性方法SLIM达到0.7161。复查使用公开的原切分和评价程序，以验证集选择训练轮数；简单基线用35次贝叶斯搜索，其中5次为初始随机配置。这些数值说明优劣随数据和基线变化，不是不同数据集的HR可以直接互比，也不是线上消费效果；第~\ref{sec:rs-evaluation-resources}~节进一步说明怎样把比较落实到信息、调参与服务成本。

搜索与广告在同一时期沿各自的问题延续。搜索从词项匹配走向学习排序、稠密和稀疏语义表示、多向量交互、标识生成与会话检索，反复面对的仍是当前需求怎样对应真实结果。广告的响应链路、支付与预算则把预测和控制连接起来；更加复杂的预估网络没有取消标签成熟、市场竞争和累计资金约束。跨域、群组、多行为、隐私与撤回也不应被挤到一条单纯的网络更替线上：跨域改变哪些实体可以对应，群组改变谁共同使用结果，多行为改变证据，隐私限制改变什么可以交换，撤回改变什么仍被允许保留。方法适应这些条件的方式可能互补，也可能冲突，应按任务及其条件比较，而不以发表年份替代解释。

归结到当前，推荐研究的主要问题仍集中在几个方向。更长历史需要解决相关访问、压缩及更新；语义与协同结合需要说明文字相近何时支持个体匹配；动态目录需要检验新物品怎样进入、失效物品怎样退出，而不是只报告静态测试集上的命中。系统还要决定怎样主动补充信息、怎样控制结果的不可靠程度：询问可以消除歧义却增加操作负担，探索可以让新物品取得证据却占用当次机会，预测不确定性能够提示证据薄弱却不能保证每次展示都正确；独立校准下的集合风险控制可以建立特定平均风险界，代价可能是输出更小的集合，并需要校准数据与待用请求具有规定的分布关系，具体条件见第~\ref{sec:rs-distribution-free-reliability}~节。连续应用进一步改变环境：旧策略决定下一轮能学到哪些反馈，持续曝光可能改变兴趣，提供者也可能调整供给，只有把数据选择、用户状态与供给响应分开建模，才知道一个闭环实验检验了什么；模拟中更好的策略仍需与真实长期效果区分。广告中的生成创意和生成竞价也遵守同一原则：创意生成提供素材候选，事实和投放资格需要另行核验，轨迹生成提供对未来状态或动作的预测，实际支出和硬预算仍由执行过程维护；追踪受限时，受保护的归因计数、聚合媒体响应曲线和随机增量实验提供不同层次的证据。模型规模、语言流畅度或统一接口，都不能替代这些应用条件。

% Outline: 1.5
\section{本章小结}
\label{sec:rs-overview-summary}

推荐、搜索与广告都需要把有限证据转为对真实对象的选择。当前需求和物品事实决定匹配依据，候选取得限定可比较的空间，分数提供特定目标下的判断，展示决定用户实际面对的内容与机会。反馈既为后续学习提供信息，也受到此前行动的影响，必须保留它的观测范围。

三者的差异体现在条件上。搜索将明确查询与相关性贯穿过程，广告增加付费主体、分配支付和跨请求预算。研究内容沿理解、召回、排序、展示，以及多方目标、长期价值和效果证据分布到后续章节；方法的历史发展可以改变模块内部计算、候选访问、共享训练或列表产生，但这些改变是否解决了应用问题，要由输出、可执行条件及相应效果证据判断。

后续从理解任务展开。当前已有的是用户表达、历史事件、情境和对象资料；还缺少的是怎样把这些证据组织成可用于本次请求的需求与匹配表示，以及在信息不足或过时时怎样修正判断。
